\section{Papangelou kernel for the Pascal point process} \label{sec:pascal}


If $X$ is a negative binomial variable with parameters $p\in (0,1)$ and $a>0$, then
\[
 (n+1)\P( X = n+1) = (1-p)^a\, \frac{(a)_{n+1}}{n!} p^{n+1} = p (a+n)\P( X = n),
\]
for all $n\in \N_0$, accordingly
\[
	\E[ X f(X)] = \E[ p(a+X) f(X+1)]
\]
for all $f\colon \N_0\to \R_+$. The following proposition gives the analogous property for Pascal point processes. For future reference we state and prove the proposition for $\sigma$-finite measures $\alpha$, allowing for configurations with infinitely many particles. Thus, let $\mathbf N$ be the space of $s$-finite counting measures on $E$. As $E$ is a Borel space, $\mathbf N$ consists of the measures $\eta$ on $E$ that are finite or countable sums of Dirac measures. The space $\mathbf N$ is equipped with the $\sigma$-algebra $\mathcal N$ generated by the counting variables $B\mapsto \eta(B)$, $B\in \mathcal E$.

A \emph{Pascal point process} with parameters $p$ and $\alpha$ is a random variable $\eta\colon(\Omega,\mathcal F, \P)\to (\mathbf N,\mathcal N)$ that satisfies the properties listed in the definition of $\rho_{p,\alpha}$: Counting variables for disjoint regions~$B_i$ are independent, and each counting variable $\eta(B)$ has a negative binomial distribution with parameters $p$ and $\alpha(B)$. For sets with infinite mass $\alpha(B) = \infty$, the counting variable $\eta(B)$ is almost surely infinite. Pascal point processes have Laplace functional
\begin{gather*}
	\E\bigl[ \e^{-\int f\dd \eta}\bigr] = \exp \biggl(\! - \int \log \biggl(\frac{1- p\,\e^{-f(x)}}{1- p}\biggr) \alpha (\dd x)\biggr) = 	\exp\Biggl(\! - \sum_{j=1}^\infty \int \bigl(1- \e^{- j f(x)}\bigr) \frac{p^j}{j} \alpha(\dd x) \Biggr)
\end{gather*}
for measurable $f \colon E \to \R_+$.

\begin{Proposition}%\label{prop:papangelou}
	Fix $p\in (0,1)$ and a $\sigma$-finite measure $\alpha$ on $E$. Let $\eta$ be a Pascal point process with parameters $p$ and $\alpha$. Then
	\[
		\E\biggl[ \int F(x,\eta)\eta(\dd x)\biggr] = \E\biggl[\int F(x,\eta+\delta_x) p (\alpha+\eta)(\dd x)\biggr]
	\]
	for all measurable $F\colon E\times \mathbf N \to \R_+$.
\end{Proposition}

The proposition says that the kernel $E\times \mathcal N\to \R_+$, $(x,B)\mapsto \kappa(x,B) = p(\alpha(B) + \eta(B))$ is a~\emph{Papangelou kernel} for the Pascal point process. For a general discussion on Papangelou kernels, we refer to \cite{MatthesWarmuthMecke} or \cite[Section 1.2.2]{Rafler2009}.

\begin{proof}
	Pascal point processes are compound Poisson; we deduce the proposition from the Mecke formula for Poisson point processes. Let $\Pi = \sum_i \delta_{(X_i,N_i)}$ be a Poisson point process on $E\times \N_0$ with intensity measure $\beta = \alpha\otimes \bigl(\sum_{n=1}^\infty \frac{p^n}{n} \delta_n\bigr)$. Then $\xi = \sum_i N_i \delta_{X_i}$ is a Pascal point process with law $\rho_{p,\alpha}$. The Mecke equation for $\Pi$ \cite[Theorem 4.1]{LastPenroseLecturesOnThePoissonProcess} and measurable $G\colon (E\times \N_0)\times \mathbf N(E\times \N_0) \to \R_+$ gives
\[
	\E\biggl[ \sum_i G( (X_i, N_i), \Pi )\biggr] = \sum_{n=1}^\infty \frac{p^n}{n} \int \E\bigl[ G\bigl( (x,n), \delta_{(x,n)} + \Pi\bigr) \bigr] \alpha (\dd x).
\]
Let $\varphi\colon E\times \mathbf N(E)\to \R_+$ be a measurable function. We apply the Mecke equation for $\Pi$ to $G\bigl((x,n),\sum_i \delta_{(x_i,n_i)}\bigr):= n \varphi\bigl(x, \sum_i n_i \delta_{x_i}\bigr)$ and get
\begin{equation}
 \label{equation: mecke1}
	\E\biggl[\sum_i N_i \varphi( X_i, \xi)\biggr] = \sum_{n=1}^\infty p^n \int \E[ \varphi(x, \xi+ n \delta_x)] \alpha(\dd x).
\end{equation}
Using \eqref{equation: mecke1} for $\varphi(x,\mu) = F(x,\mu + \delta_x)$, we obtain
\begin{align}
 p\, \E \biggl[ \sum_i N_i F\bigl( X_i, \xi +\delta_{X_i}\bigr)\biggr] &= p\sum_{n=1}^\infty p^n \int \E [ F(x, \xi + (n+1) \delta_{x}) ]\alpha(\dd x) \notag \\
 &= \sum_{n=2}^\infty p^n \int \E [ F(x, \xi + n \delta_{x}) ] \alpha(\dd x). \label{eq:mecke2}
\end{align}
We apply~\eqref{equation: mecke1} to $\varphi(x,\mu) = F(x,\mu)$ and eliminate the sum over $n\geq 2$ using~\eqref{eq:mecke2}. This gives
\begin{equation*}
	\E\biggl[\sum_i N_i \varphi( X_i, \xi)\biggr] =
		p \int \E[ F(x, \xi+ \delta_x)] \alpha(\dd x) + p\, \E \biggl[ \sum_i N_i F\bigl( X_i, \xi +\delta_{X_i}\bigr)\biggr].
\end{equation*}
Equivalently,
\[
	\E\biggl[ \int F(x,\xi) \xi(\dd x)\biggr] = p\, \E\biggl[ \int F(x, \xi+\delta_x) (\alpha + \xi)(\dd x)\biggr]
\]
and the proof is complete as $\xi$ and $\eta$ are equal in distribution.
\end{proof}

Turning back to the law $\rho_{p,\alpha}$ on $\mathbf N_{<\infty}$ for finite measures $\alpha$, we obtain
\begin{equation} \label{eq:papangelou}
	\int \biggl( \int F(x,\eta)\eta(\dd x)\biggr) \rho_{p,\alpha}(\dd x)
	= \int\biggl( \int F(x,\eta+\delta_x) p (\alpha+\eta)(\dd x)\biggr) \rho_{p,\alpha}(\dd x)
\end{equation}
for all measurable $F\colon E\times \mathbf N_{<\infty}\to \R_+$, hence also all bounded measurable functions $F$ (here we use that the expected total number of particles is finite when $\alpha$ has finite total mass).

\begin{Lemma} \label{lem:adjoints}
	For all bounded measurable $\varphi\colon E\to \C$ and all $f,g\in \mathcal D$,
	\[
		\bigl\langle f, k^+(\varphi) g\bigr\rangle = \langle k^-(\varphi) f,g \rangle.
	\]
\end{Lemma}

\begin{proof}	
	Let us denote integration with respect to $\rho_{p,\alpha}$ as $\E[f(\eta)] =\int f\dd \rho_{p,\alpha}$. We apply equation~\eqref{eq:papangelou} to $F(x,\eta) = \varphi(x) \overline{f(\eta)} g(\eta - \delta_x)$ and obtain
	\begin{align*}
		\bigl\langle f, k^+(\varphi) g \bigr\rangle
			& = \frac 1{\sqrt p} \E\biggl[\int \varphi(x) \overline{f(\eta)} g(\eta - \delta_x) \eta(\dd x)\biggr]\\
			& = \sqrt p\, \E\biggl[\int \varphi(x) \overline{f(\eta + \delta_x)} g(\eta) (\alpha+ \eta)(\dd x)\biggr]\\
			& = \langle k^-(\varphi) f,g \rangle. \qedhere
\tag*{\qed}
\end{align*}
\renewcommand{\qed}{}
\end{proof}

\section{Symmetries. Proof of Theorem~\ref{thm:symmetry}} \label{sec:symmetry}

Before we turn to the proof of Theorem~\ref{thm:symmetry}, we shows that consistency implies conservativity, which is of interest in its own.

\begin{Proposition}
 \label{proposition: consistency implies conservative}
	Every consistent Markov semigroup is conservative.
\end{Proposition}

\begin{proof}
	Let $(P_t)_{t\geq 0}$ be Markov semigroup on $(\mathbf N_{<\infty}, \mathcal N_{<\infty})$ that is consistent, i.e., $P_t \mathcal A f = \mathcal A P_t f$ for all $t\geq 0$ and all measurable $f\colon \mathbf N_{<\infty} \to \R_+$. Then, for all $k\in \N$,
	\begin{equation} \label{eq:ptak}
		P_t \mathcal A^k \one = \mathcal A^k P_t \one = \mathcal A^k \one.
	\end{equation}
	A straightforward computation shows that $\mathcal A^k\one$ is a descending factorial power of the total number of particles,
	\begin{equation} \label{eq:ak-factorial}
		\bigl( \mathcal A^k \one \bigr)(\eta) = \eta(E)( \eta(E)-1) \cdots ( \eta(E) - k +1),
	\end{equation}
	for all $\eta \in \mathbf N_{<\infty}$. Let $(\eta_t)_{t\geq 0}$ be a Markov process with semigroup $(P_t)_{t\geq 0}$ and deterministic initial condition $\eta_0=\mu \in \mathbf N_{<\infty}$, defined on some probability space $\bigl(\Omega,\mathscr F, \mathbb P_\mu\bigr)$. Equations~\eqref{eq:ptak} and~\eqref{eq:ak-factorial} show that under $\mathbb P_\mu$, at all times $t\geq 0$, the total number of particles $\eta_t(E)$ has the same factorial moments as $\eta_0(E)$. The factorial moments of $\eta_0(E) = \mu(E)$ (hence $\eta_t(E)$) vanish for $k\geq \mu(E)+1$, therefore the moment problem is uniquely solvable and we conclude $\eta_t(E) = \eta_0(E) = \mu(E)$, $\P_\mu$-almost surely, and the process is conservative.
\end{proof}

Remember that the set $\mathcal D\subset \mathcal H$ consists of the bounded measurable functions $f$ supported in $\{\eta\colon \eta(E)\leq n_f\}$ for some $n_f<\infty$.

\begin{Lemma} \label{lem:closure}
	For every $\xi \in \C$, the operator $\frac{1}{\mathrm i} \bigl( \xi k^+(\1) - \overline{\xi} k^-(\1)\bigr)$ with domain $\mathcal D$ is closable and its closure is self-adjoint.
\end{Lemma}

Put differently, the operator is \emph{essentially self-adjoint} on $\mathcal D$ and the latter is a \emph{core} for the self-adjoint closure \cite[Section VIII.2]{reedSimonI}.

\begin{proof}
	The operator $A = \frac{1}{\mathrm i} \bigl( \xi k^+(\1)- \overline{\xi} k^-(\1)\bigr)$ is similar to the Segal field operators for free quantum fields, we adapt the proof of essential self-adjointness for the latter from Reed and Simon \cite[Theorem X.41\,(a)]{reedSimonII}.
	The domain $\mathcal D$ is dense in $\mathcal H$ and satisfies $\la f, A g\ra = \la Af, g\ra$ for all $f,g\in \mathcal D$ by Lemma~\ref{lem:adjoints}. Thus, $A$ with domain $\mathcal D(A) = \mathcal D$ is symmetric. Symmetric operators are always closable. 	
	To show that the closure is self-adjoint, we check that every vector $f\in \mathcal D$ is \emph{analytic} for $A$ and conclude with Nelson's analytic vector theorem \cite[Theorem~X.39]{reedSimonII}.
 As a reminder, Nelson's analytic vector theorem states that a symmetric operator on a Hilbert space whose domain contains a total set of analytic vectors is essential self-adjoint.
 A vector~$f$ is analytic for $A$ if there exists $\eps>0$ such that $\sum_{n=0}^\infty n!^{-1} \eps^n \| A^n \varphi\|<\infty$.
	Let $f\in \mathcal D$ and $m=m_f\in \N$ be such that $|f(\eta)|\leq \|f\|_\infty \1_{\{\eta(E)\leq m_f\}}$ on $\mathbf N_{<\infty}$. Then
	\begin{align*}
		&\bigl| k^+(\1) f(\eta)\bigr| \leq \frac{1}{\sqrt p} \|f\|_\infty (m_f+1) \1_{\{\eta(E) \leq m_f+1\}},\\
		&| k^-(\1) f(\eta)| \leq \sqrt p\, \|f\|_\infty \bigl(\alpha(E)+m_f-1\bigr) \1_{\{\eta(E) \leq m_f-1\}}
	\end{align*} 	
	for all $\eta \in \mathbf N_{<\infty}$. Set $\beta:= \alpha(E)$, we obtain a common upper bound to $k^\pm(\1) f$:
	\[
		\bigl| k^\pm(\1) f(\eta)\bigr| \leq \frac{1}{\sqrt p} \|f\|_\infty (\beta +m_f+1) \1_{\{\eta(E) \leq m_f + 1\}}.
	\]
	We expand the power $\bigl(k^+(\1) - k^-(\1)\bigr)^n f$ and bound each term in the expansion by repeated use of the previous inequality. This gives
	\[
		\bigl| \bigl(k^+(\1) - k^-(\1)\bigr)^n f(\eta)\bigr| \leq \biggl(\frac{ 2 \| f\|_\infty}{\sqrt p}\biggr)^{n} ( \beta + m_f + 1)_n \1_{\{\eta(E)\leq m_f+n\}}.
	\] 	
	The indicator on the right has norm $\leq 1$. Set $s = 2|\xi|\, \|f\|_\infty/\sqrt p$, we obtain
	\[
		\sum_{n=0}^\infty \frac{\eps^n}{n!} \bigl\| \bigl(\xi k^+(\1) -\overline{\xi} k^-(\1)\bigr)^n f \bigr\|
			 \leq \sum_{n=0}^\infty \frac{(s\eps)^n}{n!} \bigl(\beta + m_f + 1)_n,
	\] 	
	which is finite for $\eps< 1/s$. Thus, every vector $f\in \mathcal D$ is analytic and $A$ is essentially self-adjoint on $\mathcal D$. 		
\end{proof}

\begin{proof}[Proof of Theorem~\ref{thm:symmetry}]
	Let $t \geq 0$. As $\rho_{p,\alpha}$ is reversible, $P_t$ is a well-defined, self-adjoint and bounded operator on $L^2(\mathbf N_{<\infty}, \mathcal N_{<\infty}, \rho_{p, \alpha})$. The operator $\exp\bigl(2 \mathrm i \theta k^0(\1)\bigr)$ acts as multiplication with $\exp( \mathrm i \theta (\alpha(E) + 2\eta(E)))$, it commutes with $P_t$ because $(P_t)_{t \geq 0}$ preserves the total number of particles (see Proposition~\ref{proposition: consistency implies conservative} below).
	
		For $\exp\bigl( \xi k^+(\1) - \overline{\xi} k^-(\1)\bigr)$, the short informal reasoning is simple: $P_t$ is self-adjoint and commutes with $k^+(\1)$, therefore it also commutes with the adjoint $k^-(\1)$ and also with the difference $\xi k^+(\1) - \overline{\xi} k^-(\1)$ and the exponential $\exp\bigl( \xi k^+(\1) - \overline{\xi} k^-(\1)\bigr)$. The precise reasoning is slightly longer because the operators are unbounded and we have to consider domains carefully. Let $f$, $g$ be non-negative functions in $\mathcal D$. Arguing as in the proof of Lemma~\ref{lem:adjoints}, we obtain $\langle k^-(\1) f, P_t g\rangle = \langle f, \mathcal A P_t g \rangle$. Therefore,
	\[
		\langle k^-(\1) f, P_t g\rangle = \langle f, \mathcal A P_t g \rangle = \langle f, P_t \mathcal A g\rangle = \bigl\langle P_t f, k^+(\1) g\bigr\rangle
	\]
	for all non-negative $f,g\in \mathcal D$. 	Taking complex conjugates and switching $f$ and $g$, we get \smash{$\bigl\langle k^+(\1) f, P_t g \bigr\rangle = \langle P_t f, k^-(\1) g\rangle $} and then, by linear combination,
	\[
		\biggl \langle \frac 1{\mathrm i} \bigl(\xi k^+(\1) - \overline{\xi} k^-(\1)\bigr) f, P_t g\biggr \rangle =
		\biggl \langle P_t f, \frac 1{\mathrm i} \bigl(\xi k^+(\1) - \overline{\xi} k^-(\1)\bigr) g\biggr \rangle
	\]
	for all $f,g\in \mathcal D$. By Lemma~\ref{lem:closure}, the operator $- \mathrm {i}\bigl(\xi k^+(\1) - \overline{\xi} k^-(\1) \bigr)$ with domain $\mathcal D$ is closable and its closure $A$ is self-adjoint. 	
	Taking limits in the above equation, we obtain
	\[
		\langle A f, P_t g \rangle = \langle P_t f, A g\rangle
	\]
	for all $f$, $g$ in the domain of $A$. In particular, for fixed $g$ in the domain of $\mathcal A$, and all $f\in \mathcal D(A)$, we have $\langle A f, P_t g\rangle = \langle f, P_t A g \rangle$, hence $P_tg$ is in the domain of $A^*$ and $A^* P_t g = P_t A g$.
	
	As $A$ is self-adjoint, we conclude that $P_t$ leaves the domain of $A$ invariant and $P_t A = A P_t$ on the domain of $A$.
It follows that $P_t$ commutes with all spectral projections of $A$ \cite[Proposition~5.26]{schmudgen2012unbounded} and then, by spectral calculus, $P_t \exp(\mathrm i A) = \exp(\mathrm i A) P_t$. 	
\end{proof}
